\section{Related work and terminology boundary}
\label{sec:related}

Paper 39 combines established ideas but assigns them a narrow, project-specific
role.  The relevant comparisons are methodological: assumption-relative
no-go reasoning, obstruction sets, trusted certificates, formal dependency
graphs, and the symbolic/group-theoretic ingredients inherited from the four
predecessors.  The literature search supporting this section is frozen in the
corrected audit snapshot; its negative result is moderate-confidence and
vocabulary-dependent.

\paragraph{No-go reasoning and obstruction theory.}
\citet{Dardashti2021NoGo} organizes no-go results by goals, frameworks,
physical assumptions, mathematical structures, and background assumptions,
and emphasizes that denying one component identifies an escape rather than
closing every research route.  Classical obstruction theory already provides
stagewise extension barriers \citep{Olum1950Obstructions}; local coefficients
are likewise established machinery \citep{KwasikSun2018Local}.  Our typed EXIT
status operationalizes the same caution in a finite audit: a changed object or
determinant category is not silently relabeled as a theorem-level failure.  We
do not claim new obstruction theory or a general no-go taxonomy.

\paragraph{Finite obstruction sets and refinement.}
Finite forbidden-obstruction characterizations are classical in graph-minor
theory \citep{RobertsonSeymour2004GraphMinorsXX}, while constructive work makes
clear that existence, enumeration, and a stopping signal are separate burdens
\citep{CattellEtAl2000ObstructionSets}.  Counterexample-guided abstraction
refinement uses failures to revise a model until a property is proved or a real
counterexample remains \citep{ClarkeEtAl2000CEGAR}.  These precedents motivate
our insistence on an explicit universe and total classifier.  They do not imply
that a historically assembled research alphabet is prospectively complete.

\paragraph{Certificates and finite mathematical graph completion.}
Proof-carrying code and foundational proof certificates separate certificate
production from checking by a trusted consumer
\citep{Necula1997PCC,HeathMiller2015ProofCertificates}.  The Equational
Theories Project is the closest functional predecessor to the broad idea of a
machine-checked finite mathematical graph \citep{BolanEtAl2025ETP}.  It fixes
4,694 short equational laws, uses formal proofs and countermodels for generating
relations, and reports the complete unrestricted implication graph.  The
project paper also records an important trust-boundary detail: transitive and
duality expansion to the full graph is performed by external programs rather
than as millions of separate end-to-end Lean theorems.  ETP's universe is
syntactically enumerable and homogeneous; ours is a retrospective,
content-addressed map of heterogeneous repairs, ownership resets, and category
exits.  The comparison rules out any claim that Paper 39 is the first exhaustive
finite mathematical graph.

\paragraph{Typed proof and dependency DAGs.}
The mathlib ecosystem provides a mature formal-library baseline
\citep{MathlibCommunity2020}.  Recent systems organize informal and formal
proof steps as dependency graphs, infer blueprint edges, track status, and use
failed nodes for refinement
\citep{CabralEtAl2025ProofFlow,ZhuEtAl2026LeanArchitect,ChungEtAl2026GoedelArchitect}.
A recent network analysis of mathlib further separates explicit,
compiler-synthesized, statement, proof, module, and namespace edges
\citep{LiEtAl2026MathlibNetwork}.  These works establish that typed DAGs and
multilayer provenance are not novel by themselves.  Paper 39 differs in what
the nodes mean: they are outcomes of authorized candidate repairs, and
completion means obstruction-or-exit coverage of a finite request contract,
not completion of one proof blueprint.

\paragraph{“Closure certificate” is an occupied term.}
\citet{MuraliEtAl2024ClosureCertificates} define closure certificates for
dynamical verification as state-pair functions that overapproximate transition
relations and establish temporal properties; the journal extension further
develops that framework \citep{MuraliEtAl2026ClosureCertificates}.  We
therefore avoid the bare label.  The phrases \emph{typed affine-program closure
audit} and \emph{retrospective obstruction-coverage audit} distinguish our
finite research-program artifact from their dynamical-systems certificate.

\paragraph{Presentation and symbolic machinery.}
Presentation-changing Tietze moves are classical \citep{Tietze1908Invarianten},
and universal decision claims over finite group presentations meet classical
undecidability boundaries \citep{Rabin1958Recursive}.  Bass's graph-of-groups
covering theory encodes tree actions \citep{Bass1993Covering}, while Krieger's
work illustrates representation-invariant symbolic objects
\citep{Krieger2000Invariant}.  Accordingly, Paper 39 calls the Paper-38 tree
canonical only for the frozen ascending-HNN splitting and presentation; it
claims no presentation-invariant decision procedure.

The constituent dynamics are also prior art: non-backtracking graph-zeta
operators \citep{Hashimoto1989Zeta}, tree-lattice Ihara determinants
\citep{Bass1992IharaSelberg}, symbolic cocycles
\citep{Schmidt1998FundamentalCocycles}, group-symbolic obstruction principles
\citep{Bartholdi2010Gardens}, Cayley-graph symbolic classifications
\citep{AubrunBitar2024SAW}, marker-based subshift embeddings
\citep{Meyerovitch2025Embedding}, and recent geometric obstructions for group
self-simulation \citep{BarbieriEtAl2025GeometricObstruction}.  Paper 39 does
not re-prove those theories.  Its qualified novelty lies only in the exact
content-addressed integration and executable relative-coverage architecture.
